\chapter*{Changelog}

\section*{EAC3}

The following changes were introduced in version EAC3.
\subsection*{AARCH-26192: Ensure validity of return types in ASL}

Return types of function and getter calls are checked after substituting the call's parameters.
For example, given the following function:
\begin{lstlisting}
func Bad{N}() => bits(N DIV 2)
begin
    var result : bits(N DIV 2);
    return result;
end;
\end{lstlisting}
the following local declaration is no longer permitted:
\begin{lstlisting}
var x = Bad{3}();
\end{lstlisting}
Its inferred type, \lstinline{bits(3 DIV 2)}, is invalid because \lstinline{DIV} requires exact division.
See \chapref{SubprogramCalls}.

Symbolic reduction no longer looks through asserting type conversions (\chapref{SymbolicEquivalenceTesting}).
The standard library function \lstinline{SInt} no longer accepts zero-width bitvectors, for example \lstinline|SInt(Zeros{0})|.

\subsection*{AARCH-26193: ASL side effects for array indices and left-hand sides}

Array index expressions in reads and writes must be readonly.
For example, if \lstinline{next_index()} is not marked \lstinline{readonly}, the following is no longer valid ASL:
\begin{lstlisting}
a[[next_index()]] = 42;
\end{lstlisting}
Except for arguments to setter calls, all subexpressions of an assignable expression must therefore be readonly.

The evaluation order of assignments is updated:
\begin{itemize}
  \item First the right-hand side is evaluated.
  \item Then array index and slice expressions on the left-hand side are evaluated exactly once, in left-to-right order.
  \item The assignment is performed with the already-evaluated array index and slice expressions.
\end{itemize}
Setter slice and field assignments retain their read-modify-write behaviour.
See \chapref{AssignableExpressions}, particularly \secref{PreparingAssignableExpressions}.

\subsection*{AARCH-26176: Clarify grammar of ASL1.0 exception catchers}

Exception catchers require a type name rather than a general type expression.
Parentheses around exception type names, such as \lstinline{when (E) => ...}, are no longer permitted.
Exception matching is unchanged.
See \chapref{CatchingExceptions}.

\subsection*{ASL-986: Source line endings and string literals}

The treatment of line endings and string literals has been clarified in the following sections:
\begin{itemize}
  \item \secref{Lexical Regular Expressions} and \secref{Whitespace}
  \item \secref{Comments}
  \item \secref{String Literals}
\end{itemize}

\section*{EAC2}

The following changes were introduced in version EAC2.

\subsection*{AARCH-25136: Improve exposition of ASL1.0's grammar construction}

The construction of ASL1.0's grammar has been clarified in the following sections:
\begin{itemize}
  \item Introduction to \chapref{Syntax}
  \item \secref{ASLGrammar}
  \item \secref{PriorityAndAssociativity}
  \item \secref{GrammarExamples}, a new section with worked examples
  \item \secref{BinaryExpressions}
\end{itemize}

\subsection*{AARCH-25137: Improve exposition of ASL1.0's evaluation order and its consequences}

Order of evaluation has been clarified in the following sections:
\begin{itemize}
  \item \secref{EvaluationOrder}, including new example \ExampleRef{Evaluation order}
  \item \secref{ParenthesisedExpressions}
\end{itemize}

\subsection*{AARCH-25350: Remove ASL1.0 shorthands for readability}

The following shorthand slice forms are no longer permitted:
\begin{itemize}
  \item Scaled form, \lstinline{factor*:len}. Equivalent to \lstinline{(factor*len) +: len}.
  \item Shorthand length form, \lstinline{:len}. Equivalent to \lstinline{0 +: len}.
\end{itemize}

Eliding of parameters on the right-hand sides of declarations with an explicit \lstinline{bits(N)} type annotation is no longer permitted.
For example, the following is no longer valid ASL:
\begin{lstlisting}
let x : bits(32) = Zeros{}();
let y : bits(64) = ZeroExtend{,1}('1');
\end{lstlisting}

\subsection*{AARCH-25640: Remove ASL1.0 constructs not used in pseudocode}

The following features are removed:
\begin{itemize}
  \item Enumeration-indexed arrays.
  \item Subtyping.
  \item Nested patterns.
\end{itemize}

\paragraph{Enumeration-indexed-arrays.}

Arrays of the form
\begin{lstlisting}
array [[enumeration_type]] of element_type
\end{lstlisting}
are no longer permitted.
For example, the following is no longer valid ASL:
\begin{lstlisting}
type Color of enumeration {RED, GREEN, BLUE};
var x : array[[Color]] of integer;
\end{lstlisting}

\paragraph{Subtyping.}

Named type declarations using the \lstinline{subtypes} or \lstinline{with} keywords are no longer permitted.
For example, the following is no longer valid ASL:
\begin{lstlisting}
type SuperInt of integer;
type SubInt of integer subtypes SuperInt;

type Packet of record { data : bits(32) };
type TimedPacket subtypes Packet with { time : integer };
\end{lstlisting}

\paragraph{Nested patterns.}

ASL patterns of the following set-like forms are no longer permitted to be nested:
\begin{lstlisting}
{pattern1, ...}
!{pattern1, ...}
\end{lstlisting}
ASL nested patterns of the following tuple form are no longer permitted:
\begin{lstlisting}
(pattern1, ...)
\end{lstlisting}
For example, the following will no longer be valid ASL.
\begin{lstlisting}
x IN {{1, 2}, 3}
y IN {!{TRUE}}
z IN {(1, 2)}
\end{lstlisting}
Note that this does not affect tuple expressions of the form:
\begin{lstlisting}
(expression1, ...)
\end{lstlisting}
